\documentclass[../DoAn.tex]{subfiles}
\begin{document}

\section{Khảo sát hiện trạng và đánh giá sản phẩm tương tự}
\label{section:2.1}
Trên thị trường hiện nay, có nhiều tựa game hành động góc nhìn thứ ba và góc nhìn thứ nhất (FPS) nổi bật tích hợp cơ chế thao tác không gian, điển hình như dòng game giải đố \textit{Portal} \cite{MetacriticPortal} hay tựa game bắn súng trực tuyến \textit{Splitgate} \cite{MetacriticSplitgate}. Tuy nhiên, phần lớn các sản phẩm này tập trung thuần túy vào yếu tố giải đố hoặc chế độ nhiều người chơi thi đấu (PvP). Thể loại bắn súng góc nhìn thứ ba kết hợp cơ chế dịch chuyển không gian theo định hướng chơi đơn (PvE - Đi theo cốt truyện) vẫn là một phân khúc còn nhiều không gian để khai thác. Paradox Caliber được thiết kế để lấp đầy khoảng trống này, kết hợp cốt truyện khoa học viễn tưởng với lối chơi hành động kịch tính. Nhân vật chính sử dụng linh hoạt khả năng mở cổng không gian để kiểm soát không gian, tạo ra các góc bắn chiến thuật nhằm hạ gục kẻ địch.

Về mặt nguồn lực và chi phí sản xuất, dự án định hướng tận dụng tối đa sức mạnh đồ họa của Unreal Engine 5.6 \cite{UE56Doc}. Thay vì thiết kế tài nguyên 3D từ đầu (gây tiêu tốn nhân lực và thời gian), dự án áp dụng chiến lược sử dụng tài nguyên (Asset) miễn phí hoặc có chi phí thấp từ nền tảng Fab/Epic Games \cite{EpicFab}. Bằng kỹ thuật tự tối ưu hóa thông số vật liệu, hiệu ứng hình ảnh (VFX) và hệ thống ánh sáng, đồ họa của trò chơi vẫn tiệm cận tiêu chuẩn của các sản phẩm thương mại chất lượng cao. Đáng chú ý, tổng mức đầu tư tài chính cho tài nguyên cấu thành dự án được kiểm soát dưới mức 1.500.000 VNĐ. Thành quả này chứng minh tính khả thi của mô hình phát triển độc lập chi phí thấp nhưng vẫn đảm bảo chất lượng kỹ thuật chuyên sâu.

\section{Mục đích của trò chơi}
\label{section:2.2}
Mục tiêu cốt lõi của dự án là cung cấp một sản phẩm giải trí tương tác cao trên nền tảng máy tính. Bên cạnh yếu tố giải trí thông thường, trò chơi đóng vai trò như một môi trường thử thách và rèn luyện phản xạ (Reflexes) ở nhịp độ cao. Thông qua cơ chế thiết lập cổng không gian đa hướng, hệ thống buộc người chơi phải vận dụng tư duy logic không gian (Spatial Logic), liên tục phân tích địa hình để vạch ra các chiến thuật (Tactical Planning) và làm chủ nhịp độ trận đấu.

\section{Định hướng sử dụng}
\label{section:2.3}
Sản phẩm được đóng gói dưới dạng một trò chơi độc lập ngoại tuyến (Offline Stand-alone Game) dành cho máy tính cá nhân. Mục đích sử dụng chính là phục vụ nhu cầu giải trí và thao tác kỹ năng của người dùng cuối. Dự án không tập trung vào việc đánh giá các định lượng tâm lý khó đo lường (như mức độ giải tỏa căng thẳng), mà định vị rõ ràng là một sản phẩm thương mại giải trí tiêu chuẩn. Đồng thời, từ góc độ học thuật, Đồ án là một công trình nghiên cứu minh chứng cho khả năng thiết kế hệ thống và tinh chỉnh đồ họa trên Unreal Engine 5.6 đối với mô hình sản xuất vi mô (Low-budget Indie).

\section{Cơ sở của định hướng sử dụng}
\label{section:2.4}
Tính khả thi và sự phù hợp trong định hướng phát triển của trò chơi được xây dựng dựa trên xu hướng thực tiễn của công nghệ và thị hiếu thị trường. Việc kết hợp sự tự do của cơ chế cổng không gian với nhịp độ hành động của TPS đã được chứng minh là yếu tố giữ chân người chơi rất tốt thông qua các cộng đồng game quốc tế.

Bên cạnh đó, chiến lược lắp ghép và tùy biến tài nguyên (Asset Integration) trên các Game Engine thế hệ mới thay vì tự sản xuất mô hình thô đang là quy chuẩn làm việc (Workflow) phổ biến của các nhà phát triển độc lập trên toàn cầu. Phương pháp này giải quyết triệt để bài toán chi phí đắt đỏ trong quy trình làm game truyền thống. Đây là cơ sở thực chứng vững chắc, đảm bảo Đồ án có thể đạt được chiều sâu kỹ thuật và chất lượng đầu ra tiệm cận với các đội ngũ phát triển chuyên nghiệp.

\end{document}